% ========== SECTION 4.4: ACRONYMS AND ABBREVIATIONS ==========

\section{Acronyms and Abbreviations}
\label{sec:acronyms}

\lettrine{T}{his section} provides a quick reference for all acronyms and abbreviations used throughout the document. Each entry includes the full form and brief context to facilitate understanding.

\subsection*{Primary Acronyms}
\addcontentsline{toc}{subsection}{Primary Acronyms}

\vspace{0.3cm}

\noindent\textbf{AIDE} --- AI-First Development Environment. A development platform architecturally designed with AI agents as primary actors.

\vspace{0.2cm}

\noindent\textbf{ADD} --- Agent-Driven Development. A paradigm where specialized AI agents collaborate to handle implementation while developers provide strategic direction.

\vspace{0.2cm}

\noindent\textbf{DEA} --- Dual Ensemble Architecture. Symphony's two-layer extension system comprising The Pit and The Grand Stage.

\vspace{0.2cm}

\noindent\textbf{IaE} --- Internal Assembling. The privileged layer of infrastructure extensions running in-process (The Pit).

\vspace{0.2cm}

\noindent\textbf{UFE} --- User-Facing Extensions. The community layer of extensions running in isolated processes (The Grand Stage).

\vspace{0.2cm}

\noindent\textbf{FQG} --- Function Quest Game. A logic puzzle system used to train the Conductor through reinforcement learning.

\vspace{0.2cm}

\noindent\textbf{ACA} --- Access Control Architecture. Symphony's three-tier trust system for extensions (Community, Trusted, Enterprise).

\subsection*{Technical Terms}
\addcontentsline{toc}{subsection}{Technical Terms}

\vspace{0.3cm}

\noindent\textbf{API} --- Application Programming Interface. A set of protocols and tools for building software applications.

\vspace{0.2cm}

\noindent\textbf{IPC} --- Inter-Process Communication. Mechanisms for exchanging data between separate processes.

\vspace{0.2cm}

\noindent\textbf{DAG} --- Directed Acyclic Graph. A graph structure used to represent dependencies and execution order.

\vspace{0.2cm}

\noindent\textbf{RL} --- Reinforcement Learning. A machine learning approach where agents learn through interaction and rewards.

\vspace{0.2cm}

\noindent\textbf{SDK} --- Software Development Kit. A collection of tools for developing extensions and applications.

\vspace{0.2cm}

\noindent\textbf{UI} --- User Interface. The visual and interactive elements through which users interact with the system.

\vspace{0.2cm}

\noindent\textbf{UX} --- User Experience. The overall experience of a person using the system.

\subsection*{System Components}
\addcontentsline{toc}{subsection}{System Components}

\vspace{0.3cm}

\noindent\textbf{IDE} --- Integrated Development Environment. A software application providing comprehensive facilities for software development.

\vspace{0.2cm}

\noindent\textbf{AI} --- Artificial Intelligence. Computer systems capable of performing tasks that typically require human intelligence.

\vspace{0.2cm}

\noindent\textbf{ML} --- Machine Learning. A subset of AI focused on systems that learn from data.

\vspace{0.2cm}

\noindent\textbf{LLM} --- Large Language Model. AI models trained on vast amounts of text data for natural language understanding and generation.

\vspace{0.2cm}

\noindent\textbf{NLP} --- Natural Language Processing. The field of AI concerned with interactions between computers and human language.

\subsection*{Development Concepts}
\addcontentsline{toc}{subsection}{Development Concepts}

\vspace{0.3cm}

\noindent\textbf{CI/CD} --- Continuous Integration and Continuous Deployment. Practices for automating software delivery.

\vspace{0.2cm}

\noindent\textbf{SLA} --- Service Level Agreement. A commitment between service provider and user defining expected service quality.

\vspace{0.2cm}

\noindent\textbf{QoS} --- Quality of Service. Measures of service performance and reliability.

\vspace{0.2cm}

\noindent\textbf{WASM} --- WebAssembly. A binary instruction format for stack-based virtual machines.

\vspace{0.2cm}

\noindent\textbf{JSON} --- JavaScript Object Notation. A lightweight data interchange format.

\vspace{0.2cm}

\noindent\textbf{YAML} --- YAML Ain't Markup Language. A human-readable data serialization format.

\subsection*{Security and Trust}
\addcontentsline{toc}{subsection}{Security and Trust}

\vspace{0.3cm}

\noindent\textbf{ACL} --- Access Control List. A list of permissions attached to an object specifying which users can access it.

\vspace{0.2cm}

\noindent\textbf{RBAC} --- Role-Based Access Control. An approach to restricting system access based on user roles.

\vspace{0.2cm}

\noindent\textbf{TLS} --- Transport Layer Security. Cryptographic protocols for secure communication.

\vspace{0.2cm}

\noindent\textbf{PKI} --- Public Key Infrastructure. A framework for managing digital certificates and public-key encryption.

\subsection*{Performance and Optimization}
\addcontentsline{toc}{subsection}{Performance and Optimization}

\vspace{0.3cm}

\noindent\textbf{CPU} --- Central Processing Unit. The primary component of a computer that performs instructions.

\vspace{0.2cm}

\noindent\textbf{RAM} --- Random Access Memory. Computer memory that can be read and changed in any order.

\vspace{0.2cm}

\noindent\textbf{I/O} --- Input/Output. Communication between an information processing system and the outside world.

\vspace{0.2cm}

\noindent\textbf{GPU} --- Graphics Processing Unit. A specialized processor designed for parallel processing tasks.

\vspace{0.2cm}

\noindent\textbf{TPU} --- Tensor Processing Unit. An AI accelerator application-specific integrated circuit.

\subsection*{Organizational Terms}
\addcontentsline{toc}{subsection}{Organizational Terms}

\vspace{0.3cm}

\noindent\textbf{BFCAI} --- Benha Faculty of Computer Science and Artificial Intelligence. The academic institution where this project was developed.

\vspace{0.2cm}

\noindent\textbf{OSS} --- Open Source Software. Software with source code that anyone can inspect, modify, and enhance.

\vspace{0.2cm}

\noindent\textbf{SaaS} --- Software as a Service. A software licensing and delivery model where software is centrally hosted.
